% Part: methods
% Chapter: proofs
% Section: introduction

\documentclass[../../../include/open-logic-section]{subfiles}

\begin{document}

\olfileid{mth}{prf}{int}

\olsection{جان پہچان}

منطق دی جان پہچان وچ اپنے تجربے کارن شاید تسی کسے !!a{derivation} نظام نوں سوکھ نال ورت لیندے او، شاید قدرتی اخذ یا فِچ ڈھنگ دا !!{derivation} نظام، یا شاید ثبوت دے رُکھاں والا نظام۔ تہانوں شاید یاد ہووے کہ ایہناں نظاماں وچ تسی ثبوت کیویں کردے سو، یا !!a{formula} ثابت کردے سو یا وکھاؤندے سو کہ دتی ہوئی دلیل جائز اے۔ ایہ کم کرن لئی تسی نظام دے قاعدے لاؤندے سو جدوں تک لوڑیندا آخری نتیجہ نہیں مل جاندا سی۔ منطق \emph{بارے} دلیل دیندیاں وی اسیں گلاں ثابت کردے آں، پر بہتے ویلے اسیں کوئی !!a{derivation} نظام نہیں ورتدے۔ اصل وچ، اسیں جیہڑے ثبوت ویکھدے آں، اوہناں وچوں بہتے انگریزی وچ دتے جاندے نیں (شاید کجھ علامتی زبان وی وچ پا دتی جائے)، سارے دے سارے پہلے درجے دی منطق دی زبان وچ نہیں۔ ایہو جیہے ثبوت بناندیاں پہلاں شاید تہانوں سمجھ نہ آوے۔ میں !!a{derivation} نظام توں بناں کوئی گل کیویں ثابت کراں؟ میں شروع کتھوں کراں؟ مینوں کیویں پتا لگے کہ میرا ثبوت ٹھیک اے؟

ثبوت کرن دی کوشش توں پہلاں ایہ جاننا ضروری اے کہ ثبوت کیہ اے تے کیویں بنائی دا اے۔ جیویں ناں توں ای پتا لگدا اے، \emph{ثبوت} دا کم وکھاؤنا اے کہ کوئی گل سچ اے۔ تسی ایس نوں اک گل بات وانگ سوچ سکدے او۔ کوئی تہانوں پُچھدا اے کہ کوئی گل سچ اے، مثلاً کی دو توں چھڈ کے ہر اوّلی عدد طاق اے۔ جواب وچ «ہاں» آکھنا کافی نہیں۔ اوہ جاننا چاہ سکدا اے کہ \emph{کیوں}۔ ایس صورت وچ تسی اوہنوں اک ثبوت دیو گے۔

روز دی گل بات وچ جواب ول اشارہ کر دینا یا ادھورا جواب دینا کافی ہو سکدا اے۔ پر منطق تے ریاضی وچ سانوں پکا ثبوت چاہیدا اے۔ اسیں وکھاؤنا چاہندے آں کہ گل \emph{کسے وی} شک توں پرے سچ اے۔ ایس دا مطلب اے کہ ساڈے ثبوت دا ہر قدم جائز ہونا چاہیدا اے، تے اوہدا جواز ٹھوس ہووے۔ یعنی جیہڑی گل تسی منّ کے ورتدے او، اوہ قضیے دے بیان وچ واقعی منّی گئی ہووے، تسی تعریفاں ٹھیک طرح لاؤ، تے جیہڑیاں دلیلاں نوں جواز بناؤندے او اوہ ٹھیک استنباط ہون، وغیرہ۔

عام طور اُتے اسیں کوئی بیان ثابت کردے آں۔ جیہڑے بیان ثابت کردے آں اوہناں دے کئی ناں نیں، قضیہ، وڈا قضیہ، معاون قضیہ، یا نتیجہ۔ قضیہ اک بنیادی، ثابت کرن جوگا بیان اے۔ ایہ درج کرن جوگا اے، پر شاید خاص ڈونگھا نہ ہووے یا بہت واری ورتیا نہ جاوے۔ وڈا قضیہ اک خاص اہم قضیہ اے۔ اوہدا ثبوت اکثر کئی قدماں وچ ونڈیا ہُندا اے، تے کدی اوہدا ناں اوہنوں پہلاں ثابت کرن والے دے ناں اُتے ہُندا اے (جیویں کانتور دا قضیہ، یا لووَن ہائم–سکولم دا قضیہ)، یا جیہڑی گل اوہ کردا اے اوہ دے ناں اُتے (جیویں تکمیل دا قضیہ)۔ معاون قضیہ اوہ قضیہ اے جیہڑا کسے ہور اہم نتیجے دے ثبوت وچ ورتیا جاندا اے۔ پر ایس گل نال بھلیکھا پے سکدا اے۔ کدی معاون قضیے اپنے آپ وچ اہم نتیجے ہُندے نیں تے پیش کرن والے دے ناں اُتے وی جاندے نیں (جیویں زورن دا معاون قضیہ)۔ نتیجہ اوہ گل اے جیہڑی کسے ہور نتیجے توں سوکھ نال نکلدی اے۔

جیہڑا بیان ثابت کرنا ہووے، اوہدے وچ اکثر منّیاں ہوئیاں شرطاں ہُندیاں نیں جیہڑیاں صاف کردیاں نیں کہ اسیں کیہڑی قسم دیاں شےواں بارے کوئی گل ثابت کردے آں۔ اوہ ایویں شروع ہو سکدا اے، «$!A$، $!B \lif !C$ دی صورت دا !!a{formula} ہووے» یا «منّ لو $\Gamma \Proves !A$»، یا ایہو جیہی کوئی گل۔ ایہ قضیے، وڈے قضیے، یا معاون قضیے دیاں \emph{مفروضہ شرطاں} نیں، تے تسی اپنے ثبوت وچ ایہناں نوں سچ منّ سکدے او۔ ایہ حد لاؤندیاں نیں کہ اسیں کی ثابت کردے آں، تے جنہاں شےواں دی گل کردے آں اوہناں دے کجھ ناں وی دیندیاں نیں۔ مثلاً جے تہاڈا قضیہ ایویں شروع ہُندا اے، «$!A$، $!B \lif !C$ دی صورت دا !!a{formula} ہووے»، تاں تسی بس اک خاص قسم دے سارے فارمولیاں (یعنی شرطیاں) بارے کوئی گل ثابت کردے او۔ تے ایہ سمجھیا جاندا اے کہ $!B \lif !C$ کوئی وی شرطیہ اے جیہدی گل تہاڈا ثبوت کرے گا۔

\end{document}
